%% LyX 2.4.4 created this file.  For more info, see https://www.lyx.org/.
%% Do not edit unless you really know what you are doing.
\documentclass[english]{article}
\usepackage[T1]{fontenc}
\usepackage[latin9]{inputenc}
\usepackage{enumitem}
\usepackage{amsmath}
\usepackage{amsthm}
\usepackage{cancel}
\usepackage{geometry}
\geometry{verbose,tmargin=1in,bmargin=1in,lmargin=1in,rmargin=1in}

\makeatletter

%%%%%%%%%%%%%%%%%%%%%%%%%%%%%% LyX specific LaTeX commands.
%% Because html converters don't know tabularnewline
\providecommand{\tabularnewline}{\\}

%%%%%%%%%%%%%%%%%%%%%%%%%%%%%% Textclass specific LaTeX commands.
\numberwithin{equation}{section}
\numberwithin{figure}{section}
\newlength{\lyxlabelwidth}      % auxiliary length 

%%%%%%%%%%%%%%%%%%%%%%%%%%%%%% User specified LaTeX commands.
\usepackage{fancyhdr}
\usepackage{lastpage}
\pagestyle{fancy}
\usepackage{enumitem}
\usepackage{ifthen}
%sets math fonts to be more readable
\everymath=\expandafter{\the\everymath\displaystyle}
%\DeclareMathSizes{10}{10}{10}{10}
\usepackage{multicol}

\usepackage{tikz}
\usetikzlibrary{plotmarks}
\usepackage{pgfplots}
\pgfplotsset{compat=newest}

\usepackage{calc}
\usepackage[nomessages]{fp}

\usepackage{advdate}    % Advancing/saving dates
\usepackage{datetime}   % Dates formatting
\usepackage{datenumber} % Counters for dates
% Added by lyx2lyx
\setlength{\parskip}{\smallskipamount}
\setlength{\parindent}{0pt}

\makeatother

\usepackage{babel}
\begin{document}
\renewcommand{\author}{Dr.\,James Ashe}

\newcommand{\classNum}{132}

\newcommand{\classSec}{K}

\newcommand{\nameType}{Name:}

\newcommand{\examType}{Exam 1}

\renewcommand{\date}{\formatdate{6}{2}{2014}}

%%First page's header%%%%%%%%%%%%%%%%%%%%%%
\fancypagestyle{firststyle} %%header settings for first pagr
{
	\fancyhead[L]{MTH \classNum{}(\classSec{}) \author{}\\
	\textbf{\examType{}} \date{}}
	\fancyhead[C]{\textbf{\nameType}}
	\setlength{\headsep}{14pt}
}
%%%%%%%%%%%%%%%%%%%%%%%%%%%%%%%%%%%%%%%%%%

%%Next page's header%%%%%%%%%%%%%%%%%%%%%%%
\setlist{nolistsep}
\lhead{MTH \classNum{}(\classSec{})} 
\chead{\textbf{\nameType}} 
\cfoot{\thepage{}~of~\pageref{LastPage}} 
\setlength{\headsep}{5pt} 
%%%%%%%%%%%%%%%%%%%%%%%%%%%%%%%%%%%%%%%%%%%

%%Various settings; see the preamble for other settings%%%%%
%%\setenumerate[0]{leftmargin=12pt,itemindent=0pt,labelwidth=0pt}
\renewcommand{\labelenumi}{\textbf{\arabic{enumi}.}}
\renewcommand{\labelenumii}{\textbf{\alph{enumii}.}}
\thispagestyle{firststyle}
%%%%%%%%%%%%%%%%%%%%%%%%%%%%%%%%%%%%%%%%%%
\newcommand{\xmax}{10}
\newcommand{\xmin}{-10}
\newcommand{\xstep}{0} %must be xmin+n
\newcommand{\ymax}{10}
\newcommand{\ymin}{-10}
\newcommand{\ystep}{0} %must be ymin+n

\newcommand{\xspace}{\vspace{1.2cm}}

\emph{You must show work to receive credit. Each problem is worth
10 pts unless stated otherwise.}

\textbf{Use the given conditions to write an equation for the line
in slope-intercept form, if possible.}
\begin{enumerate}[resume]
\item Slope$=5$, passing through $(1,3)$.\vfill
\item Through $(2,-3)$ and $(-3,4)$.\vfill
\end{enumerate}
\textbf{Write a cost function for the problem. Assume that the relationship
is linear.}
\begin{enumerate}[resume]
\item Alfred owns a small publishing company. Printing a book has a fixed
cost of \$500 plus \$3 per book printed. Find the cost function for
printing $x$ number of books. \xspace
\end{enumerate}
\textbf{Write a revenue function for the problem.}
\begin{enumerate}[resume]
\item If Alfred's publishing company sells books for \$12 a piece, find
the revenue function for selling $x$ number of books.\xspace
\end{enumerate}
\textbf{Solve the problem.}
\begin{enumerate}[resume]
\item \label{enu:Producing--units}Producing $x$ units of tacos costs $C(x)=5x+20$;
revenue is $R(x)=15x$, where $C(x)$ and $R(x)$ are in dollars.

\begin{enumerate}[resume]
\item Find the profit function for selling $x$ number of tacos.\vfill
\item How many tacos must be sold to make a profit of \$120?\vfill
\end{enumerate}
\end{enumerate}
\newpage
\begin{enumerate}[resume]
\item Suppose that the supply and demand functions for honey is 
\[
p=S(q)=0.3q+2.7\mbox{ and }p=D(q)=6.9-0.4q
\]
where $p$ is the price in dollars for an 8-oz container and $q$
is the number of barrels.

\begin{enumerate}[resume]
\item Find the equilibrium quantity.\vfill
\item Find the equilibrium price. \vfill
\end{enumerate}
\end{enumerate}
\textbf{Find the equation of the least squares line. }
\begin{enumerate}[resume]
\item \label{enu:The-total-amount}The total amount of consumer credit has
been increasing in recent years. The following tables gives the total
US outstanding consumer credit in billions of dollars for selected
years.

\begin{enumerate}[resume]
\item []%
\begin{tabular}{c|c|c|c|c|c}
Year & 2004 & 2005 & 2006 & 2007 & 2008\tabularnewline
\hline 
Consumer Credit & 2219.5 & 2319.5 & 2415.0 & 2551.9 & 2592.1\tabularnewline
\end{tabular}\xspace\xspace
\end{enumerate}
\end{enumerate}
\textbf{SHORT ANSWER. Write the word or phrase that best completes
each statement or answers the question.}
\begin{enumerate}[resume]
\item {[}5 points{]} Explain what it means to ``break even'' in problem
\ref{enu:Producing--units}.\xspace\xspace
\item {[}5 points{]} The correlation coefficient from problem \ref{enu:The-total-amount}
is $r\approx0.991$. Does this indicate a good fit of the least squares?
If another data set has a correlation coefficient of $r=-.999$, does
the data fit more or less closely to a line than the data from problem
\ref{enu:The-total-amount}? \xspace\xspace
\end{enumerate}

\end{document}
